\begin{example}
    La sucesión $(an^\sigma)_{n\geq1}$ con $a \in \mathbb{R}$, $a \neq 0$ y 
    $\sigma \in (0,1)$ es equidistribuida módulo $1$.
\end{example}

\begin{resolve}    
    Para demostrarlo utilizaremos el criterio de Weyl. 

    Sea $k \in \mathbb{Z}\setminus\{0\}$, como $a \neq 0$, entonces $b = ak \neq 0$. 
    Basta comprobar que 
    \[
        \frac{1}{N} \sum_{n=1}^N e^{2\pi i b n^\sigma} \longrightarrow 0 
        \qquad \text{cuando} \ N \to \infty.
    \]

    Aplicando la fórmula (\ref{prop:formula_euler}) a la función $f(x) = e^{2\pi i b x^\sigma}$,
    \[
        \sum_{n=1}^N e^{2\pi i \, b n^\sigma}
        = \int_1^N e^{2\pi i b x^\sigma} \, dx 
        + \frac{1}{2} \left(e^{2\pi i \, b N^\sigma} + e^{2\pi i b}\right)
        + \int_1^N \left(\partfrac{x} - \frac{1}{2}\right) 2\pi i b \sigma x^{\sigma-1} \, e^{2\pi i b x^\sigma} \, dx 
    \]

    Calculamos la primera integral usando integración por partes,
    \begin{align*}
        \int_1^N e^{2\pi i b x^\sigma} \, dx 
        = \int_1^N \frac{\left(e^{2\pi i b x^\sigma}\right)'}{2\pi i b \sigma x^{\sigma-1}} \, dx
        &= \frac{1}{2 \pi i b \sigma} \int_1^N \left(e^{2\pi i b x^\sigma}\right)'x^{1-\sigma} \, dx \\
        &= \frac{1}{2 \pi i b \sigma} \left(\left[e^{2\pi i b x^\sigma}x^{1-\sigma}\right]_1^N + (1-\sigma) \int_1^N e^{2\pi i b x^\sigma} x^{-\sigma} \, dx \right)
    \end{align*}  
    
    Acotando cada uno de los términos anteriores,
    \begin{align*}
        \left| e^{2\pi i b N^\sigma}N^{1-\sigma} - e^{2\pi i b}\right|
        &\leq N^{1-\sigma} + 1 \\
        \left|\int_1^N e^{2\pi i b x^\sigma} x^{-\sigma} \, dx\right|
        &\leq \int_1^N x^{-\sigma} \, dx
        = \frac{N^{1-\sigma} - 1}{1-\sigma}
    \end{align*}
    Ambos crecen del orden de $N^{1-\sigma}$, por lo que la primera integral es$O(N^{1-\sigma})$.

    El término central se acota por
    \[
        \left|\frac{1}{2} \left(e^{2\pi i b N^\sigma} + e^{2\pi i b}\right)\right| \leq 1
    \]
    luego es $O(1)$.

    Para la segunda integral, usando que $|\partfrac{x} - \frac{1}{2}| \leq \frac{1}{2}$,
    y que $|e^{2\pi i b x^\sigma}| = 1$, se tiene
    \[
        \left|\int_1^N \left(\partfrac{x} - \frac{1}{2}\right) 2\pi i b \sigma x^{\sigma-1} \, e^{2\pi i b x^\sigma} \, dx\right|
        \leq \pi |b| \int_1^N \sigma x^{\sigma-1} \, dx  
        = \pi |b| (N^\sigma - 1)
    \]
    por lo que esta integral es $O(N^\sigma)$.
    
    Reuniendo las estimaciones anteriores, y teniendo en cuenta que 
    $\sigma \in (0,1)$, se obtiene
    \[
        \sum_{n=1}^N e^{2\pi i \, b n^\sigma} = O(N^{1-\sigma}) + O(1) + O(N^\sigma) = O(N^{\max\{\sigma, 1-\sigma\}}) = o(N)
    \]
    De este modo,
    \[
        \frac{1}{N} \sum_{n=1}^N e^{2\pi i b n^\sigma} \longrightarrow 0 
        \qquad \text{cuando} \ N \to \infty,
    \]
    y por tanto la sucesión $(an^\sigma)$ es equidistribuida módulo $1$.
\end{resolve}

